\clearpage
\phantomsection% 
\addcontentsline{toc}{chapter}{RINGKASAN}
\thispagestyle{fancy}

\begin{center}
	\large \bfseries \MakeUppercase{Ringkasan}\\
	\normalsize \normalfont {\thetitle}\\
	\normalsize \normalfont {\theauthor}\\
	\bigskip
	
	\normalsize \normalfont \justifying
	\normalsize \normalfont \justifying
Ketidakseimbangan kelas merupakan permasalahan utama dalam klasifikasi \textit{machine learning}, di mana model cenderung bias terhadap kelas mayoritas. Eroglu dan Pir (2024) mengusulkan HOUM yang menggabungkan Safe-Level SMOTE dan SVM secara iteratif, namun belum mengeksplorasi variasi teknik \textit{oversampling} serta konfigurasi SVM. Oleh karena itu, penelitian ini bertujuan menganalisis pengaruh variasi teknik \textit{oversampling} Safe-Level SMOTE dan ADASYN serta konfigurasi SVM (parameter $C$ dan jenis kernel) pada HOUM terhadap kinerja klasifikasi data tidak seimbang.

Penelitian menggunakan empat dataset \textit{benchmark} yaitu Pima Indians Diabetes (rasio 1:1,87), Haberman \textit{Survival} (1:2,78), \textit{Blood Transfusion} (1:3,20), dan Glass5 (1:22,78) yang ditambahkan untuk menguji HOUM pada ketidakseimbangan ekstrem. Data dibagi 80:20 menggunakan \textit{stratified split}. HOUM diimplementasikan dengan 18 kombinasi konfigurasi SVM dengan nilai parameter $C$ 0,5; 1; 1,5; 2; 2,5; 3 dan jenis kernel Linear, RBF, Polynomial yang dikombinasikan dengan Safe-Level SMOTE dan ADASYN. Penyeimbangan dilakukan iteratif hingga rasio kelas seimbang. Data hasil penyeimbangan melatih \textit{Random Forest} dan dievaluasi menggunakan \textit{Recall}, \textit{Precision}, G-Mean, serta ROC-AUC.

Hasil penelitian menunjukkan HOUM-SLSMOTE unggul pada \textit{Recall} di Diabetes, Haberman, dan Transfusion, sedangkan HOUM-ADASYN unggul pada \textit{Precision} di Diabetes dan Glass5. Pada Glass5, HOUM-ADASYN mencapai \textit{Recall} 1,0000, ROC-AUC 1,0000, dan G-Mean 0,9877. Nilai $C$ kecil (0,5) optimal untuk Diabetes, Haberman, dan Glass5, sementara $C=2$ optimal untuk Transfusion. RBF unggul pada data berdimensi tinggi (Diabetes), Linear lebih stabil pada data berdimensi rendah (Haberman, Transfusion), sedangkan pada Glass5 seluruh kernel menghasilkan metrik identik pada HOUM-ADASYN. G-Mean tertinggi 0,9877 dicapai pada Glass5 dengan HOUM-ADASYN, diikuti Diabetes 0,7401. Peningkatan \textit{Recall} signifikan di seluruh dataset meskipun disertai penurunan \textit{Precision} pada beberapa kasus.

Penelitian menyimpulkan bahwa efektivitas teknik \textit{oversampling} bergantung pada karakteristik dataset, pemilihan parameter SVM harus disesuaikan dimensi data, dan HOUM efektif meningkatkan deteksi kelas minoritas termasuk pada kasus ketidakseimbangan ekstrem. Direkomendasikan eksplorasi dataset berdimensi lebih tinggi, rentang $C$ di bawah 0,5, teknik \textit{oversampling} alternatif, serta algoritma klasifikasi lainnya.

\vfill

\end{center}
\clearpage